\section{Conclusion}
\label{sec:ch3_conclusion}

We introduced \textbf{AERO}, a lightweight and adaptive forecasting model for resource management in edge-cloud networks. AERO dynamically detects and adjusts to dominant periodicities in workload data, overcoming the limitations of fixed-periodicity models. With its minimal parameter count, it is suitable for resource-constrained environments.

Evaluations on real-world datasets and live deployments demonstrate that \ac{AERO} achieves comparable prediction accuracy and operational performance to models with highly advanced architectures and millions of parameters, such as Pathformer, while greatly reducing complexity and computational overhead. Simulations using the extended \ac{COSCO} framework show that integrating \ac{AERO} enhances energy efficiency, reduces response times, and minimizes \ac{SLA} violations compared to reactive approaches. Our live deployment experiments further validate these findings, demonstrating that \ac{AERO} delivers strong performance.

Adaptive prediction models effectively manage fluctuating workloads in dynamic edge-cloud environments. Future work includes integrating \ac{AERO} into real-time orchestration systems and examining its scalability in larger deployments.